% !TEX root = ../main.tex
\section{Hạn chế và Thách thức}

Mặc dù hệ thống đã chứng minh được hiệu quả trong việc giải quyết bài toán xếp lịch ở quy mô hiện tại, quá trình thực nghiệm cũng chỉ ra một số thách thức và hạn chế cần được khắc phục trong các giai đoạn phát triển tiếp theo. Một trong những vấn đề trọng tâm là tính khả mở (scalability) của hệ thống. Khi số lượng lớp học phần tăng lên đến hàng trăm hoặc hàng nghìn, không gian tìm kiếm sẽ bùng nổ vượt quá khả năng xử lý hiệu quả của thuật toán di truyền truyền thống. Trong những trường hợp này, việc duy trì tính đa dạng của quần thể mà vẫn đảm bảo tốc độ hội tụ đòi hỏi các cấu trúc dữ liệu phức tạp hơn và có thể cần đến sự hỗ trợ của các kỹ thuật lập trình song song hoặc phân tán.

Bên cạnh đó, mô hình ràng buộc hiện tại vẫn chưa bao quát được toàn bộ các khía cạnh đa dạng của thực tiễn giáo dục đại học. Các yếu tố mang tính định tính như sở thích cá nhân của giảng viên về khung giờ dạy, khoảng cách địa lý giữa các tòa nhà trong khuôn viên trường, hay các yêu cầu đặc thù về trang thiết bị của từng môn học cần được tích hợp sâu hơn vào hàm đánh giá. Việc mở rộng hệ thống ràng buộc này sẽ làm tăng độ phức tạp của hàm thích nghi và đòi hỏi các toán tử đột biến có định hướng rõ ràng hơn để tránh việc phá vỡ các cấu trúc lịch đã được tối ưu một phần.

Từ góc độ thuật toán, sự phụ thuộc vào các giá trị khởi tạo ngẫu nhiên và các tham số cố định như tỷ lệ lai ghép hay đột biến cũng là một điểm cần cải thiện. Một hệ thống lý tưởng nên có khả năng tự điều chỉnh các thông số này dựa trên trạng thái tiến hóa của quần thể (adaptive GA) nhằm tránh hiện tượng hội tụ sớm vào các cực trị địa phương. Việc kết hợp thuật toán di truyền với các phương pháp tìm kiếm cục bộ (Local Search) như Hill Climbing hay Tabu Search cũng là một hướng đi hứa hẹn nhằm tinh chỉnh chất lượng lịch học và giảm thiểu các chỉ số phạt mềm một cách triệt để hơn. Những nỗ lực này không chỉ giúp nâng cao tính thực tiễn của hệ thống mà còn mở ra cơ hội áp dụng rộng rãi hơn trong các bài toán tối ưu hóa tài nguyên có quy mô lớn.
